% Generated by report/make_restated.py from theory/A1_lower_bound.tex: do not edit (edit the note and rerun).
\begin{restated}{Proposition}{prop:A4}{A1}{1}{testable prediction of hypothesis A4}
Fix $n\ge2$, $m\ge2$, $\mu$ and a relative accuracy $\varepsilon\in(0,\frac1{105})$, start at $\theta_0=0$, and let $\kappa$ vary. On the instance \eqref{eq:instance} with $\zeta=\zeta_Q$ ($Q$, $q$ as in Proposition~\ref{prop:minimizer}; $\mu_F$, $L_F$ as in Proposition~\ref{prop:constants}), let $T_{\rm alg}(\varepsilon)$ be the number of gradient evaluations (its expectation for L-Katyusha) that an algorithm needs to reach $\E\norm{\theta-\etastar}^2\le\varepsilon\norm{\theta_0-\etastar}^2$: SAGA with $\gamma=\frac1{4R^2}$ and $z_i^{(0)}=\fip(0)$, the $n$ gradients of this sweep being counted; L-Katyusha with the parameters of Theorem~11 of \cite{kovalev2020dont}, $\theta$ being its mirror sequence $z^k$, charged $n$ gradients for $\Fp(w^0)$ and, per iteration, one gradient plus $n$ when the reference point is refreshed (probability $p=\frac1n$; the gradients $\fip(w^k)$ of each full pass are stored), i.e.\ $2$ per iteration in expectation; GD and AGD on $F$, $n$ per iteration. Let $C_0$ be the constant of \eqref{eq:LKconst}.
\begin{enumerate}[nosep,label=(\roman*)]
\item (Plateau, $\kappa\le n$.) For SAGA, under \eqref{ineq:dim} and $\varepsilon<q^2/2$,
\[
n+(n-1)\log\frac{q^2}{2\varepsilon}\ \le\ T_{\rm SAGA}(\varepsilon)\ \le\ n+\max\Big\{3n,\frac{16\kappa}3\Big\}\log\frac{1+n/4}\varepsilon\ \le\ n+\frac{16n}3\log\frac{1+n/4}\varepsilon ,
\]
both sides being $\Theta(n\log(1/\varepsilon))$ once $\varepsilon\ll q^2$ and $\varepsilon\ll1/n$ (the upper bound is $\Theta(n\log(n/\varepsilon))$ in general): the plateau with slope $\approx0$ is a theorem for SAGA. For L-Katyusha only
\[
\frac{n\log(1/(2\varepsilon))}{2\log(1/q)}\ \le\ T_{\rm LK}(\varepsilon)\ \le\ n+16n\log\frac{C_0}\varepsilon,\qquad C_0\le n+\tfrac73,
\]
is proved (the lower bound under \eqref{ineq:dim}); the two sides are of the same order $n\log(1/\varepsilon)$ only when $\log(1/q)=O(1)$, i.e.\ when $\kappa$ is of order $n$ (for $\kappa=n$, $\log(1/q)=2\operatorname{artanh}\big((2-\frac1n)^{-1/2}\big)$, which decreases to $\log(3+2\sqrt2)\approx1.763$ as $n\to\infty$: $1.81$ for $n=16$, $1.77$ for $n=100$). For $\kappa\ll n$, $\log(1/q)\approx\log\frac{4n}{\kappa-1}$, and a gain of this logarithmic order over $n\log(1/\varepsilon)$ is not excluded for L-Katyusha; it is realized by GD on $F$ with step $1/L_F$, which contracts the error by $1-\frac{\mu_F}{L_F}\le1-\frac1Q=\frac{\kappa-1}{n+\kappa-1}\approx\frac{\kappa-1}n$ per full gradient, in agreement with Theorem~\ref{thm:HLOY}(c) and with item (ii) of the assessment of hypothesis (A1) in Section~\ref{sec:lb:statements}. In the upper bounds, the second term takes over at $\kappa=\frac9{16}n$ for SAGA ($\frac{16\kappa}3=3n$) and at $\kappa=\frac83n$ for L-Katyusha (its momentum parameter $\theta_1$ reaches $\frac12$).
\item (Power laws.) Let
\begin{equation*}
\tag{\ref*{eq:window}}
\kappa_{\max}(m,\varepsilon):=1+n\min\Big\{\frac{16m^2}{\log^2(1/\varepsilon)}-1,\ \frac{(m+1)^2}{\pi^2}\Big\},
\end{equation*}
where the first term is the sufficient form $m\ge\frac{\sqrt Q}4\log\frac1\varepsilon$ of the dimension condition \eqref{ineq:dim} solved for $\kappa$ (so that $q^{2m}\le\varepsilon$), and the second is the condition $y\ge2$ of Lemma~\ref{lem:overlap} solved for $\kappa$. In the window $4n+1\le\kappa\le\kappa_{\max}(m,\varepsilon)$, with $\rho=\rho(m,\kappa)$ the slow-mode fraction of Lemma~\ref{lem:overlap},
\begin{align*}
\tag{\ref*{eq:sagawindow}}
\Big(\kappa-\frac12\Big)\log\frac{\rho}{\varepsilon}\ &\le\ T_{\rm SAGA}(\varepsilon)\ \le\ n+\frac{16\kappa}3\log\frac{1+n/4}\varepsilon,\qquad \rho=\rho(m,\kappa)\ge\frac1{105},\\
\tag{\ref*{eq:lkwindow}}
\frac{\sqrt{n(\kappa-1)}}8\log\frac1{2\varepsilon}\ &\le\ T_{\rm LK}(\varepsilon)\ \le\ n+4\sqrt6\,\sqrt{n\kappa}\,\log\frac{C_0}\varepsilon,\qquad C_0\le1+\frac58Q ,
\end{align*}
so that $T_{\rm SAGA}=\Theta(\kappa\log(1/\varepsilon))$ and $T_{\rm LK}=\Theta(\sqrt{n\kappa}\log(1/\varepsilon))$ up to the displayed constants and the $\log(1+\frac58Q)$ of \eqref{eq:lkwindow}: $\kappa$ enters $T_{\rm SAGA}$ and $T_{\rm LK}$ with the exponents $1$ and $\frac12$ as orders of growth (not as the slope of a line fitted over the window, see Remark~\ref{rem:plot}), and $T_{\rm SAGA}/T_{\rm LK}\asymp\sqrt{\kappa/n}$. GD on $F$ (step $1/L_F$) and AGD on $F$ have the same two orders of growth: $n(\frac{L_F}{4\mu}-\frac12)\log\frac\rho\varepsilon\le T_{\rm GD}\le\frac{nQ}2\log\frac1\varepsilon$ with $L_F\ge\mu(1+\frac{13}{20}(Q-1))$, and $\frac{\sqrt{n(\kappa-1)}}8\log\frac1{2\varepsilon}\le T_{\rm AGD}=O(n\sqrt Q\log\frac1\varepsilon)$.
\item (Saturation.) For $\kappa\gg nm^2$ the chain is exhausted: $\mu_F/\mu=1+\frac{\kappa-1}{4n}\lmin(A_m(\zeta_Q))$ grows linearly in $\kappa$, so
\begin{equation*}
\tag{\ref*{eq:kappaeff}}
\kappa_{\rm eff}=\frac{R^2}{\mu_F}=\frac{\kappa}{1+\frac{\kappa-1}{4n}\lmin(A_m(\zeta_Q))}\ \xrightarrow[\kappa\to\infty]{}\ \frac{4n}{\lmin(A_m(1))}=\frac{n}{\sin^2\frac{\pi}{2(2m+1)}}\approx\frac{4n(2m+1)^2}{\pi^2},
\end{equation*}
because $\zeta_Q\to1$ and $\lmin$ is continuous in the matrix; for every finite $\kappa$, Lemma~\ref{lem:chain} gives the two-sided bound $\frac{\kappa}{1+\frac{\kappa-1}{n}\sin^2\frac{\pi}{2(m+1)}}\le\kappa_{\rm eff}\le\frac{\kappa}{1+\frac{\kappa-1}{n}\sin^2\frac{\pi}{2(2m+1)}}$. Then $(2\kappa_{\rm eff}-\frac12)\log\frac\delta\varepsilon\le T_{\rm SAGA}(\varepsilon)\le n+\frac{16\kappa_{\rm eff}}3\log\frac{1+n/4}\varepsilon$ for $\varepsilon<\delta$, with the single-mode overlap $\delta=\inner{x^\ast}{v_1}^2/\norm{x^\ast}^2$ ($x^\ast$, $v_1$ as in Lemma~\ref{lem:overlap}), which tends to $\delta_\infty(m)>0$ (Remark~\ref{rem:rho}): $T_{\rm SAGA}$ saturates at $\Theta(\kappa_{\rm eff}\log(1/\varepsilon))=\Theta(nm^2\log(1/\varepsilon))$, while $T_{\rm LK}$ is no longer bounded below by $\sqrt{n\kappa}$ (the hypothesis $q^{2m}\le\varepsilon$ of Proposition~\ref{prop:anyseq} fails) and its behaviour depends on the strong-convexity constant used to set the parameter $\sigma$ of L-Katyusha. Since $F$ is $\mu_F$-strongly convex, Theorem~11 of \cite{kovalev2020dont} also applies with $\sigma=\mu_F/R^2$ (the tuning used in the experiments, where $\mu=\lmin(\nabla^2F)$ is passed to the method): \eqref{eq:LKconst}--\eqref{eq:LKupper} then hold with $\mu_F$ and $\kappa_{\rm eff}$ in place of $\mu$ and $\kappa$, so that $T_{\rm LK}=O\big((n+\sqrt{n\kappa_{\rm eff}})\log\frac{C_0}\varepsilon\big)$ with $C_0$ bounded in $\kappa$, and $T_{\rm LK}$ saturates at $O(nm\log\frac{C_0}\varepsilon)$. With the nominal $\sigma=\mu/R^2$ of \eqref{eq:LKupper} no saturation is proved: the upper bound keeps growing like $\sqrt{n\kappa}$, and the momentum $\theta_1=\sqrt{2n\mu/(3R^2)}$ tends to $0$, whereas the value matched to the curvature, $\sqrt{2n\mu_F/(3R^2)}$, does not. Thus $T_{\rm SAGA}$, and $T_{\rm LK}$ with $\sigma=\mu_F/R^2$, stay bounded as $\kappa\to\infty$, and their log-log slopes against $\kappa$ tend to $0$.
\end{enumerate}
\end{restated}
